% Part: turing-machines
% Chapter: machines-computations
% Section: turing-machines

\documentclass[../../../include/open-logic-section]{subfiles}

\begin{document}

\olfileid{tur}{mac}{tur}
\olsection{ټيورينګ ماشينونه}

\begin{explain}
د ټيورينګ ماشين د جوړښت صوري تعريف مجرد ښکاري، خو په اصل کښې ساده
دے: دا يوازې هغه ټول معلومات په يوه رياضياتي جوړښت کښې رانغاړي چې د
ټيورينګ ماشين د کار د ټاکلو دپاره پکار دي۔ په دې کښې دا شامل دي: (1)
ماشين په کومو حالتونو کښې اوسېدلے شي، (2) پر پټه کومې نښې روا دي، (3)
ماشين بايد په کوم حالت کښې پيل وکړي، او (4) د ماشين د لارښوونو سټ څه
دے۔
\end{explain}

\begin{defn}[ټيورينګ ماشين]
يو \emph{ټيورينګ ماشين} $M$ هغه څلور‌يز $\langle Q, \Sigma, q_0,
\delta\rangle$ دے چې له لاندې توکو جوړ دے:
\begin{enumerate}
\item د \emph{حالتونو} يو متناهي سټ~$Q$،
\item يوه متناهي \emph{الفبا} $\Sigma$ چې $\TMendtape$ او
  $\TMblank$ پکښې شامل دي،
\item يو \emph{لومړنی حالت}~$q_0 \in Q$،
\item د \emph{لارښوونو} يو متناهي سټ~$\delta\colon Q \times \Sigma
  \pto Q \times \Sigma \times \{\TMleft, \TMright, \TMstay\}$۔
\end{enumerate}
جزوي تابع~$\delta$ ته د~$M$ \emph{د حالت بدلون تابع} هم ويل کېږي۔
\end{defn}

\begin{explain}
موږ فرض کوو چې پټه يوازې په يوه لوري کښې نامتناهي ده۔ ځکه نو ګټوره
ده چې ځانګړې نښه~$\TMendtape$ د پټې د کيڼ پای د نښې په توګه وټاکو۔
په دې سره د ټيورينګ ماشين پروګرامونه په اسانۍ پوهېدے شي چې کله له پټې
څخه د ``لوېدو په خطر'' کښې دي۔ موږ دا فرض کولے شو چې دا نښه هېڅکله
نۀ بدلېږي، يعنې که $\delta(q,\TMendtape)$ تعريف شوې وي، نو
$\delta(q,\TMendtape) = \tuple{q', \TMendtape, x}$ وي۔ ځينې درسي
کتابونه دا کار کوي؛ موږ ئې نۀ کوو۔ د خپل ټيورينګ ماشين د جوړولو پر
وخت يوازې دومره پام وکړئ چې~$\TMendtape$ هېڅکله بدله نۀ کړي۔ سربېره
پر دې، داسې حالتونه هم شته چې د دې نښې د بدلولو اجازه څه ګټور
انعطاف برابروي۔
\end{explain}

\begin{ex}
\emph{جفت ماشين:} جفت ماشين په صوري ډول هغه څلور‌يز
$\tuple{Q, \Sigma, q_0, \delta}$ دے چې
\begin{align*}
Q & = \{ q_0, q_1 \} \\
\Sigma & = \{ \TMendtape, \TMblank, \TMstroke \}, \\
\delta(q_0, \TMstroke) & = \tuple{q_1, \TMstroke, \TMright},\\
\delta(q_1, \TMstroke) & = \tuple{q_0, \TMstroke, \TMright},\\
\delta(q_1, \TMblank)  & = \tuple{q_1, \TMblank, \TMright}.
\end{align*}
\end{ex}

\end{document}
